\documentclass[11pt]{article}
\usepackage{mathblatt}
\begin{document}
\blattkopf{Konfidenzintervalle}{Lernblatt}
\verzeichniszeile{\verz{e1}{1 Das Intervall lesen und deuten: die Überdeckungsdeutung (verträglich heißt: der angenommene Anteil wird vom Intervall überdeckt (Nr.~9–10)} \verztrenn \verz{e2}{2 Mit der Näherungsformel rechnen: die Grenzgleichung nach p lösen (quadratisch (Nr.~11–13)} \verztrenn \verz{e3}{3 Argumente über Formel und Verträglichkeit: die Länge des Intervalls bei doppeltem Umfang (Faktor eins durch Wurzel zwei (Nr.~14–15)} \verztrenn \verz{abhaken}{Das kann ich}}
% Einheit 1 von 3 – Das Intervall lesen und deuten: die Überdeckungsdeutung (verträglich heißt: der angenommene Anteil wird vom Intervall überdeckt (bank/konfidenzintervalle/e1.jsonl, zusammenbau v0.1)
%% TODO Zeitmarke Einheit 1: Marken-Zeile nicht lesbar
%% TODO Zweigzeile Teil 1: was hier gelernt wird (Satz in Schülersprache aus dem Einheitstitel) – Einheit 1
\einheitenkopf[e1]{Einheit 1 von 3 · Das Intervall lesen und deuten: die Überdeckungsdeutung (verträglich heißt: der angenommene Anteil wird vom Intervall überdeckt}
\zweigzeile{keine Prüfungsaufgabe}
\setcounter{aufgabe}{8}

%% TODO Titel: Ich-kann-Satz fehlt in der Bank (Platzhalter: Kettenname) – Nr. 9
%% TODO Anweisung: ein Satz über der ersten Teilaufgabe fehlt in der Bank – Nr. 9
\begin{aufgabe}{Lesen und Überdeckung}
%% TODO \rechenplatz{n}: Schrittzahl des Grundfalls fehlt in der Bank (nur bei fester Schreibform, 2.3 b)
\begin{teile}
\teil Richtig oder falsch gedeutet? Nur ankreuzen: „Der Anteil $0{,}3$ ist mit dem Ergebnis verträglich, weil das Konfidenzintervall ihn überdeckt.“ \\ \kreuz{richtig gedeutet} \\ \kreuz{falsch gedeutet}
\teil Die Grafik zeigt Grenzgraphen für Stichproben vom Umfang $n = 100$ (Sicherheitswahrscheinlichkeit 95\,\%). In einer Stichprobe gibt es 30 Treffer. Lies das Konfidenzintervall ab.

\begin{ksys}[xmin=0,xmax=1,ymin=0,ymax=1,xstep=0.1,ystep=0.1,xlabel=p,ylabel=h,ablesen] \funktion{\x-1.96*sqrt(\x*(1-\x)/100)}{g_1} \funktion{\x+1.96*sqrt(\x*(1-\x)/100)}{g_2} \end{ksys}
[ \leerfeld ; \leerfeld ]
\teil Vermutet wird ein Anteil von 50\,\%. In einer Stichprobe vom Umfang $n = 100$ gibt es 42 Treffer. Lies das Konfidenzintervall zur Sicherheitswahrscheinlichkeit 95\,\% ab und beurteile die Vermutung.

\begin{ksys}[xmin=0,xmax=1,ymin=0,ymax=1,xstep=0.1,ystep=0.1,xlabel=p,ylabel=h,ablesen] \funktion{\x-1.96*sqrt(\x*(1-\x)/100)}{g_1} \funktion{\x+1.96*sqrt(\x*(1-\x)/100)}{g_2} \end{ksys}
\teil Die Grafik zeigt die Konfidenzintervalle aus 8 Stichproben (Anteile in Prozent). Lies ab: Gib einen Anteil an, der mit genau 6 der 8 Stichprobenergebnisse verträglich ist.

\begin{zahlengerade}[xmin=30,xmax=65] \intervall[1]{38}{52}{1} \intervall[2]{42}{56}{2} \intervall[3]{35}{49}{3} \intervall[4]{44}{58}{4} \intervall[5]{40}{54}{5} \intervall[6]{47}{61}{6} \intervall[7]{36}{50}{7} \intervall[8]{43}{57}{8} \end{zahlengerade}
\leerfeld \%
\teil Aus 25 weiteren Stichproben vom Umfang 400 wird jeweils das Konfidenzintervall zur Sicherheitswahrscheinlichkeit 90\,\% bestimmt; der wahre Anteil p bleibt gleich. Y ist die Anzahl der Intervalle, die p überdecken. Gib die Verteilung von Y mit Begründung an und weise nach: Die Wahrscheinlichkeit, dass genau 23 Intervalle p überdecken, ist kleiner als 27\,\% \hfill (Abitur 2026 LK).
\end{teile}
\end{aufgabe}

%% TODO Titel: Ich-kann-Satz fehlt in der Bank (Platzhalter: Kettenname) – Nr. 10
%% TODO Anweisung: ein Satz über der ersten Teilaufgabe fehlt in der Bank – Nr. 10
\begin{aufgabe}{Lesen und Überdeckung – Fehler finden, Begründen}
\begin{teile}
\teil Die Grafik zeigt die Grenzgraphen zur Sicherheitswahrscheinlichkeit 95\,\% für $n = 50$. Bei 40 Treffern schneidet Mia die Senkrechte bei $p = 0{,}8$ mit beiden Graphen und liest $[0{,}69; 0{,}91]$ ab. Finde den Fehler und lies richtig ab.

\begin{ksys}[xmin=0,xmax=1,ymin=0,ymax=1,xstep=0.1,ystep=0.1,xlabel=p,ylabel=h,ablesen] \funktion{\x-1.96*sqrt(\x*(1-\x)/50)}{g_1} \funktion{\x+1.96*sqrt(\x*(1-\x)/50)}{g_2} \end{ksys}
\teil Begründe, warum die Anzahl der Konfidenzintervalle, die den wahren Anteil überdecken, binomialverteilt ist, und nenne ihre Trefferwahrscheinlichkeit.
\end{teile}
\end{aufgabe}

\clearpage
% Einheit 2 von 3 – Mit der Näherungsformel rechnen: die Grenzgleichung nach p lösen (quadratisch (bank/konfidenzintervalle/e2.jsonl, zusammenbau v0.1)
%% TODO Zeitmarke Einheit 2: Marken-Zeile nicht lesbar
%% TODO Zweigzeile Teil 1: was hier gelernt wird (Satz in Schülersprache aus dem Einheitstitel) – Einheit 2
\einheitenkopf[e2]{Einheit 2 von 3 · Mit der Näherungsformel rechnen: die Grenzgleichung nach p lösen (quadratisch}
\zweigzeile{keine Prüfungsaufgabe}
\setcounter{aufgabe}{10}

%% TODO Titel: Ich-kann-Satz fehlt in der Bank (Platzhalter: Kettenname) – Nr. 11
%% TODO Anweisung: ein Satz über der ersten Teilaufgabe fehlt in der Bank – Nr. 11
\begin{aufgabe}{Näherungsformel}
%% TODO \rechenplatz{n}: Schrittzahl des Grundfalls fehlt in der Bank (nur bei fester Schreibform, 2.3 b)
\begin{teile}
\teil Vorwärts oder rückwärts? Nur ankreuzen: „Von 500 Befragten sagen 140 ja. Bestimme das Konfidenzintervall.“ \\ \kreuz{vorwärts: Grenzen gesucht} \\ \kreuz{rückwärts: Ergebnis oder Umfang gesucht}
\end{teile}
\begin{gleichungsraster}[2]
\gl{\text{In einer Stichprobe vom Umfang 400 gibt es 120 Treffer. Berechne mit der Grenzgleichung die Grenzen des Konfidenzintervalls zur Sicherheitswahrscheinlichkeit 95\,\% auf Tausendstel.}} &
\gl{\text{Die untere Grenze eines Konfidenzintervalls zur Sicherheitswahrscheinlichkeit 95\,\% ist $0{,}319$, der Umfang 500. Bestimme die Trefferzahl der Stichprobe.}} \\
\gl{\text{Bei einem Glücksrad nimmt man an, dass das Feld „Stern“ mit der Wahrscheinlichkeit $0{,}2$ kommt. Beobachtet wird, dass genau 25\,\% der Drehungen „Stern“ ergeben. Ermittle die kleinste Anzahl von Drehungen, bei der dieser Anteil zur Sicherheitswahrscheinlichkeit 95\,\% nicht mit $0{,}2$ verträglich ist (Abitur 2023 LK).}} & \\
\end{gleichungsraster}
\end{aufgabe}

%% TODO Titel: Ich-kann-Satz fehlt in der Bank (Platzhalter: Kettenname) – Nr. 12
%% TODO Anweisung: ein Satz über der ersten Teilaufgabe fehlt in der Bank – Nr. 12
\begin{aufgabe}{Verträglichkeit einer vermuteten Trefferwahrscheinlichkeit über das 1,96σ-Intervall um den Erwartungswert prüfen}
\begin{teile}
\teil Für Blumenzwiebeln wird eine Austriebswahrscheinlichkeit von 70\,\% vermutet; von 80 gesetzten Zwiebeln treiben 50 aus. Die Vermutung gilt zur Sicherheitswahrscheinlichkeit 95\,\% als verträglich, wenn 50 im Intervall von $\mu - 1{,}96\sigma$ bis $\mu + 1{,}96\sigma$ liegt. Untersuche das \hfill (Abitur 2017 LK).
\end{teile}
\end{aufgabe}

%% TODO Titel: Ich-kann-Satz fehlt in der Bank (Platzhalter: Kettenname) – Nr. 13
%% TODO Anweisung: ein Satz über der ersten Teilaufgabe fehlt in der Bank – Nr. 13
\begin{aufgabe}{Näherungsformel – Fehler finden, Begründen}
\begin{teile}
\teil Zu 90 Treffern bei Umfang 300 soll das Konfidenzintervall zu 95\,\% bestimmt werden. Jan rechnet $0{,}3 \pm 1{,}96 \cdot \sqrt{0{,}3 \cdot 0{,}7 / 300}$ und erhält $[0{,}248; 0{,}352]$. Finde den Fehler und rechne richtig.
\teil Begründe, warum die Grenzgleichung zwei Lösungen für p hat und warum man beide braucht.
\end{teile}
\end{aufgabe}

\clearpage
% Einheit 3 von 3 – Argumente über Formel und Verträglichkeit: die Länge des Intervalls bei doppeltem Umfang (Faktor eins durch Wurzel zwei (bank/konfidenzintervalle/e3.jsonl, zusammenbau v0.1)
%% TODO Zeitmarke Einheit 3: Marken-Zeile nicht lesbar
%% TODO Zweigzeile Teil 1: was hier gelernt wird (Satz in Schülersprache aus dem Einheitstitel) – Einheit 3
\einheitenkopf[e3]{Einheit 3 von 3 · Argumente über Formel und Verträglichkeit: die Länge des Intervalls bei doppeltem Umfang (Faktor eins durch Wurzel zwei}
\zweigzeile{keine Prüfungsaufgabe}
\setcounter{aufgabe}{13}

%% TODO Titel: Ich-kann-Satz fehlt in der Bank (Platzhalter: Kettenname) – Nr. 14
%% TODO Anweisung: ein Satz über der ersten Teilaufgabe fehlt in der Bank – Nr. 14
\begin{aufgabe}{Argumente}
%% TODO \rechenplatz{n}: Schrittzahl des Grundfalls fehlt in der Bank (nur bei fester Schreibform, 2.3 b)
\begin{teile}
\teil Stimmt die Aussage? Nur ankreuzen: „Bei vierfachem Umfang und gleichem Stichprobenanteil wird das Konfidenzintervall etwa halb so lang.“ \\ \kreuz{stimmt} \\ \kreuz{stimmt nicht}
\teil Der Stichprobenanteil ist $0{,}35$, die Annahme $p = 0{,}3$. Begründe ohne Rechnung, dass die obere Grenze des Konfidenzintervalls größer als die Annahme ist.
\teil Ein Konfidenzintervall ist bei Umfang 400 etwa $0{,}08$ lang. Wie lang ist es etwa bei gleichem Stichprobenanteil und Umfang 1600? etwa \leerfeld
\teil Der Stichprobenanteil ist $0{,}32$, die Sicherheitswahrscheinlichkeit 95\,\%. Aussage: „Obwohl $0{,}32$ genau in der Mitte zwischen $0{,}28$ und $0{,}36$ liegt, kann $0{,}28$ unverträglich und $0{,}36$ verträglich sein.“ Rechnung I: $0{,}32 = 0{,}28 + 1{,}96 \cdot \sqrt{0{,}28 \cdot 0{,}72 / n}$ ergibt $n \approx 484{,}0$; Rechnung II: $0{,}32 = 0{,}36 - 1{,}96 \cdot \sqrt{0{,}36 \cdot 0{,}64 / n}$ ergibt $n \approx 553{,}2$. Beurteile die Aussage mit beiden Rechnungen \hfill (Abitur 2025 LK).
\end{teile}
\end{aufgabe}

%% TODO Titel: Ich-kann-Satz fehlt in der Bank (Platzhalter: Kettenname) – Nr. 15
%% TODO Anweisung: ein Satz über der ersten Teilaufgabe fehlt in der Bank – Nr. 15
\begin{aufgabe}{Argumente – Fehler finden, Begründen}
\begin{teile}
\teil Bei Umfang 500 ist ein Konfidenzintervall $0{,}1$ lang. Tim sagt: „Bei Umfang 1000 und gleichem Anteil ist es $0{,}05$ lang.“ Finde den Fehler und gib die Länge richtig an.
\teil Begründe, warum der Stichprobenanteil immer im Konfidenzintervall liegt.
\end{teile}
\end{aufgabe}

% Abhakseite „Das kann ich“ – Zone nur im Gesamt (\mitzone)
%% TODO Ich-kann-Sätze: Platzhalter wie in den Aufgabendateien
\begin{abhakseite}
\ifdefined\mitzone
\abhakgruppe{Kennst du schon}
\abhak{1}{Relative Häufigkeit und Anteil aus einer Stichprobe bilden}
\abhak{2}{Binomialverteilung ansetzen und Punktwahrscheinlichkeiten berechnen}
\abhak{3}{Graphen zweier Funktionen ablesen (Schnitt mit Waagerechten und Senkrechten)}
\abhak{4}{Die Sigma-Regeln und ihre c-Werte}
\abhak{5}{Quadratische Gleichungen und Wurzelgleichungen lösen, Lösungen deuten}
\abhak{6}{Ungleichungen in n lösen und runden}
\abhak{7}{Binomialverteilung ansetzen und Punktwahrscheinlichkeiten berechnen – Fehler finden}
\abhak{8}{Binomialverteilung ansetzen und Punktwahrscheinlichkeiten berechnen – selbst rechnen}
\fi
\abhakgruppe{Einheit 1 · Das Intervall lesen und deuten: die Überdeckungsdeutung (verträglich heißt: der angenommene Anteil wird vom Intervall überdeckt}
\abhak{9}{Lesen und Überdeckung}
\abhak{10}{Lesen und Überdeckung – Fehler finden, Begründen}
\abhakgruppe{Einheit 2 · Mit der Näherungsformel rechnen: die Grenzgleichung nach p lösen (quadratisch}
\abhak{11}{Näherungsformel}
\abhak{12}{Verträglichkeit einer vermuteten Trefferwahrscheinlichkeit über das 1,96σ-Intervall um den Erwartungswert prüfen}
\abhak{13}{Näherungsformel – Fehler finden, Begründen}
\abhakgruppe{Einheit 3 · Argumente über Formel und Verträglichkeit: die Länge des Intervalls bei doppeltem Umfang (Faktor eins durch Wurzel zwei}
\abhak{14}{Argumente}
\abhak{15}{Argumente – Fehler finden, Begründen}
\end{abhakseite}

\end{document}
